\documentclass{article}
\input{../../../../lib.sty}

\title{\texttt{NoisePrivacyMap<L1Distance<RBig>, MultiDP> for ZExpFamily<1>}}
\author{\MichaelShoemateAuthor}
\date{}

\begin{document}
\maketitle

\contrib
This proves the privacy-map implementation for the discrete Laplace mechanism
when the output measure is \rustdoc{measures/struct}{MultiDP}.

\section{Hoare Triple}
\subsection*{Precondition}
\texttt{MI} is \rustdoc{metrics/type}{L1Distance<RBig>} and \texttt{MO} is
\rustdoc{measures/struct}{MultiDP}. The distribution scale is a
non-negative rational.

\subsection*{Pseudocode}
\lstinputlisting[language=Python,firstline=2,escapechar=|]{./pseudocode/NoisePrivacyMap_for_ZExpFamily1_MultiDP.py}

\subsection*{Postcondition}
\begin{theorem}
For every input distance accepted by the returned privacy map, the map returns
a \texttt{PrivacyGuarantee} containing valid privacy representations for adding an
independent discrete Laplace sample to each coordinate.
\end{theorem}

\begin{proof}
Lines \ref{line:neg-scale} and \ref{line:neg-sens} reject invalid
parameters. With zero sensitivity, line \ref{line:zero-sens} returns the
identity guarantee. With zero scale and positive sensitivity, line \ref{line:zero-scale}
returns the empty guarantee, which imposes no privacy constraint.

For positive rational sensitivity $s$ and scale $b$, the discrete Laplace
point-mass ratio calculation gives the pure-DP bound
\[
    \varepsilon = s/b.
\]
The upward cast makes the floating-point \texttt{epsilon} an upper bound on
this value, so the $(\varepsilon,0)$ point inserted into the guarantee is valid.
To justify using the scalar bounds for a vector and a rational sensitivity,
write $a=1/b$ and let $d\geq 0$ denote a shift magnitude. For integer $d$,
the scalar RDP expression evaluated by the implementation is
\[
   g_\alpha(d)=\frac{\log(A e^{(\alpha-1)ad}+B e^{-\alpha ad})}{\alpha-1},
\]
where, with $k=(2\alpha-1)a$,
\[
 A=\tanh(a/2)\left(\frac{1}{e^a-1}+\frac{1}{1-e^{-k}}\right),
 \qquad
 B=\tanh(a/2)\left(\frac{1}{e^a-1}-\frac{1}{e^k-1}\right).
\]
Both coefficients are positive and sum to one. This formula extends to real
$d\geq0$ as a convex log-sum-exp function with $g_\alpha(0)=0$.
Its right derivative at zero is nonnegative: indeed,
\[
 B=\tfrac12\left(1-\frac{\coth(k/2)}{\coth(a/2)}\right)
 \leq \frac{\alpha-1}{2\alpha-1},
\]
because $x\coth(x/2)$ is increasing for positive $x$ (its derivative
has the sign of $\sinh(x)-x$). Thus $g_\alpha$ is nondecreasing.

Likewise, the scalar discrete-Laplace zCDP bound used in the implementation
extends to $h(d)=ad-a(1-e^{-ad})/\sinh(a)$, which is convex and
nondecreasing on $d\geq 0$ and satisfies $h(0)=0$.

For either function $g$, convexity and $g(0)=0$ imply
$g(d_i)\leq (d_i/D)g(D)$ when $D=\sum_i d_i>0$.
The shifts $d_i$ are integers, and independence gives additive Rényi
divergence (and additive zCDP bounds). Hence
\[
    \sum_i g(d_i)\leq g(D)\leq g(s)
\]
whenever $D\leq s$, including rational $s$. The case $D=0$ is trivial.
This establishes the validity of the scalar formulas at the supplied
sensitivity for the vector mechanism.

The zCDP and RDP calculations retain the exact rational sensitivity, scale,
and ratio, enclosing each rational with outward-rounded endpoints before
using directed interval arithmetic. Therefore the returned $\rho$ and RDP
values are no smaller than their exact bounds. Finally,
\texttt{rdp\_discrete\_laplace} evaluates the exact Rényi expression with
directed intervals. If that evaluation fails,
the implementation first uses the exact RDP value implied by pure DP and, if
those numerics also fail, uses the elementary pure-DP bound
$D_\alpha \leq \varepsilon$. Each result is an upper bound at every valid
order. Each representation is therefore simultaneously valid, and querying or
composing the resulting \texttt{PrivacyGuarantee} preserves conservativeness.
\end{proof}

\end{document}
